\begin{table}[ht]
\centering
\caption{Ajuste dos tempos medidos à complexidade teórica.}
\label{tab:ajustes}
\small
\begin{tabular}{llrrrrr}
\hline
Algoritmo & f(n,m) & c & R² (c·f) & a (ms/vértice) & b (ms/aresta) & R² (a·n+b·m) \\
\hline
BFS & O(n+m) & 0,000088 & 0,9910 & -0,000876 & 0,000144 & 0,9952 \\
DFS & O(n+m) & 0,000090 & 0,9764 & -0,001197 & 0,000166 & 0,9832 \\
Euler (verificação) & O(n+m) & 0,000101 & 0,9767 & -0,001327 & 0,000184 & 0,9834 \\
Dijkstra & O((n+m) log n) & 0,000022 & 0,9978 & -0,002712 & 0,000455 & 0,9960 \\
Bellman-Ford & O(n·m) [efetivo O(k·m)] & 0,000054 & 0,9968 & -0,000703 & 0,000919 & 0,9970 \\
Tarjan (pontes/articulações) & O(n+m) & 0,000194 & 0,9948 & -0,001335 & 0,000284 & 0,9970 \\
Prim & O(m log n) & 0,000102 & 0,9810 & -0,033018 & 0,003261 & 0,9898 \\
Kruskal & O(m log m) & 0,000035 & 0,9963 & -0,004277 & 0,000855 & 0,9949 \\
\hline
\end{tabular}
\end{table}
